% Generated draft: requires mathematical and source review.
\documentclass[11pt]{article}
\usepackage[margin=25mm]{geometry}
\usepackage{fontspec}
\setmainfont{Latin Modern Roman}
\usepackage{amsmath,amssymb,mathtools,unicode-math}
\setmathfont{Latin Modern Math}
\usepackage{enumitem,needspace,xurl,hyperref}
\setlength{\parindent}{0pt}
\setlength{\parskip}{5pt}
\setlength{\emergencystretch}{3em}
\newcommand{\N}{\mathbb{N}}
\newcommand{\Z}{\mathbb{Z}}
\newcommand{\Q}{\mathbb{Q}}
\newcommand{\R}{\mathbb{R}}
\newcommand{\C}{\mathbb{C}}
\newcommand{\F}{\mathbb{F}}
\newcommand{\A}{\mathbb{A}}
\newcommand{\PP}{\mathbb{P}}
\newcommand{\E}{\mathbb{E}}
\newcommand{\eps}{\varepsilon}
\newcommand{\dd}{\,\mathrm{d}}
\begin{document}
\section*{UnsolvedMath ID: 30004377}
\textbf{Source problem code:} catalogue entry
\subsection{Definitions and mathematical statement}
Let \(C\) denote the Cantor set referred to in the supplied source, and let \((C_j)\) denote its stated standard prefractal sequence. The record does not specify the construction or ambient space of \(C\) and \(C_j\). Let \(\phi_j^h\) denote the source's boundary-element approximations and \(\phi\) the intended limiting solution; their defining equations and function spaces are not supplied. Impose the stated parameter choices
\[
 i(j)=j,\qquad \delta=0.
\]
The requested convergence assertion is
\[
 \phi_j^h\longrightarrow\phi\qquad\text{as }j\to\infty.
\]
This is an unresolved formulation of the source's requested assertion: the record specifies neither the topology of convergence nor whether \(h\) is fixed, depends on \(j\), or tends to zero under an additional condition. It also supplies no governing boundary-value or integral equation, boundary conditions, admissible data, or discrete approximation spaces. Consequently, the displayed assertion cannot be made into a fully quantified mathematical statement from the supplied record alone.
\subsection{Short English statement}
The problem asks whether the boundary-element approximations associated with the standard Cantor prefractals converge to the intended limiting solution when \(i(j)=j\) and \(\delta=0\). The defining problem, approximation scheme, mesh limit, and meaning of convergence require clarification before a precise theorem can be stated.
\subsection{Sources}
\begin{itemize}
\item Supplied local UnsolvedMath registry record, permanent ID \(30004377\), source problem code \texttt{catalogue entry}, in the Oberwolfach Reports Open Problems collection. No individual proposer is identified; the supplied proposed year is 2020.
\item Supplied report citation: \emph{Boundary Element Methods (2021)}, Oberwolfach Reports, DOI 10.4171/owr/2020/5, pp.\ 273--376. \url{https://doi.org/10.4171/owr/2020/5}. The background attributes the problem to workshop open-problem discussions documented in this report.
\item Supplied literature source: \emph{Boundary element methods for fractal screens}. \url{https://link.springer.com/article/10.1007/s00211-021-01182-y}.
\item Literature assessment checked at \texttt{2026-08-22}: the supplied curation labels the problem open and reports convergence results for several refined mesh and prefractal regimes, while stating that the cited literature does not prove convergence for the standard Cantor prefractal sequence with \(i(j)=j\) and \(\delta=0\). This is supplied dated curation, not an independently verified literature assessment.
\item The registry labels the statement \texttt{corrected\_verified} and reports that the official report was checked and its notation repaired. These are supplied provenance claims; no independent source verification is asserted here.
\end{itemize}
\end{document}
